\documentclass[11pt,a4paper]{article}
\usepackage[utf8]{inputenc}
\usepackage{geometry}
\geometry{margin=1in}
\usepackage{listings}
\usepackage{syntax}
\usepackage{forest}
\usepackage{booktabs}
\usepackage{hyperref}

\title{Design Specification: PARD Programming Language}
\author{Erdinç Daşkın}
\date{\today}

\begin{document}
\maketitle
\tableofcontents
\newpage

\section{Introduction \& Philosophy}
PARD is a programming language designed around game design and development.
It leverages FORTRAN's fast math and C interoperability.
The name PARD is an abbreviation of Punching Card, which is an homage to the early days of FORTRAN.
\section{Lexical Grammar}
\subsection{Variables and their types}
Variables, just like most other programming languages, are a significant part of PARD. \\
There's 4 basic variable types:
\begin{itemize}
    \item \textbf{int}: A standard integer type for counts, loop counters, and IDs.
    \item \textbf{float}: A floating-point number for positions, physics calculations, and deltatime.
    \item \textbf{bool}: A boolean value (\texttt{true} or \texttt{false}) for state flags.
    \item \textbf{string[N]}: A fixed-length, ultra-fast character array of size $N$.
\end{itemize}
As you can see, strings are not dynamically sized - this is not only for easy operation in fortran, but it also introduces a good performance boost.
\subsection{Variable definition}
You can easily define a variable in the following way: \\
\texttt{int myInt = 10} \\
\texttt{float myFloat = 10.5} \\
\texttt{bool myBool = true} \\ 
\texttt{string[6] myString = 'Hello!'}

\section{Compiler Architecture}

\end{document}